\documentclass[11pt]{article}
\usepackage[margin=1in]{geometry}
\usepackage[T1]{fontenc}
\usepackage{lmodern}
\usepackage{xcolor}
\usepackage{enumitem}
\usepackage[hidelinks]{hyperref}

\newcommand{\pkg}[1]{\textsf{#1}}
\newcommand{\code}[1]{\texttt{#1}}
\newenvironment{editorcomment}
  {\par\smallskip\sloppy\noindent\textbf{Editor comment.}\itshape}
  {\par\smallskip\fussy}
\newenvironment{response}
  {\noindent\textbf{Response.}}
  {\par\smallskip}

\title{Response to the Editor\\
\large IntegMultiReg: An R Package for Integrative Bayesian
Multi-Regression of Multi-Platform Biomarkers}
\author{Sinian Zhang, Jianfeng Wang, and Thierry Chekouo}
\date{Revision draft: September 30, 2026}

\begin{document}
\maketitle

We thank the editor for the constructive assessment. We have revised the
software interface, fit inspection, uncertainty reporting and replication
materials. The current submission uses package version 0.2.0. Below we respond
point by point; comments are condensed for readability. References to sections
use their titles so that they remain valid after manuscript pagination changes.

\section*{1. Software design and validated data objects}
\begin{editorcomment}
Further software development is needed beyond packaging the original study
code. The input interface is complicated and potentially error prone; a
dedicated data type should support validation, printing, summaries, component
extraction and prediction data.
\end{editorcomment}
\begin{response}
We introduced \code{imr\_data()}, a reusable multi-platform data class shared
by fitting and prediction. It validates subject identifiers, unique column
names, finite numerical values, outcome coding and cross-platform alignment,
and records observed-platform patterns, eligible subjects and subgroup sizes.
It supports \code{print()}, \code{summary()}, \code{as.data.frame()},
\code{validate\_imr\_data()} and standard named-component extraction. The
component-wise fitting interface remains available and passes through the
same validation path. The revised manuscript's package-interface section
explains the class and uses it in the main example.

The boundary audit also corrected integer overflow in computational controls,
ambiguous column names and incomplete validation of saved fits.
\end{response}

\section*{2. Fit inspection, uncertainty and covariates}
\begin{editorcomment}
Additional methods are needed to inspect, validate and compare fits, examine
other parameter traces, obtain posterior uncertainty and consider suitable
covariates.
\end{editorcomment}
\begin{response}
The package now supplies \code{validate\_imr()}, \code{compare\_imr()},
selection-indicator and theta-interaction traces, \code{posterior\_summary()}
and \code{confint()}. These complement the original log-posterior trace and
point summaries. \code{posterior\_draws()} provides coefficient uncertainty
and posterior prediction for continuous, binary and right-censored outcomes;
conditional sampling resamples latent responses. The methods table and the
fitted-object and diagnostic sections describe their use. Tests exercise
these methods for all three outcome families, including prediction routing,
response scales and reproducibility.

The scope is explicit: model weights use the original Laplace approximation;
conditional split R-hat does not diagnose convergence of the selection chain;
and \code{compare\_imr()} itself is a descriptive summary.
Log-time survival predictions use time-scale medians because the corresponding
posterior mean need not exist. Clinical covariates remain forced within each fitted specification, while the
selection prior applies to platform biomarkers. To respond to predictive
covariate choice without changing that prior, we added a complete worked
comparison of two prespecified formulas on the same subjects and identical CV
folds. We also added nested CV: inner folds choose the formula and outer folds
evaluate that choice. The script saves the actual subject/fold assignments and
checks that outer test IDs never enter inner validation. It is included in the
replication bundle and as an installed package example.

On the synthetic example, the two candidates have paired mean MSEs of 0.349793
and 0.376508; the separately evaluated selection procedure has outer MSE
0.340617. These are illustrative predictive results, not clinical evidence.
Necessary adjustment variables must remain in every candidate formula. This
provides an executable covariate-choice workflow; it does not implement a new
clinical inclusion prior or automate causal confounder selection. The new
``Comparing and choosing covariate specifications'' subsection explains the
settings and interpretation.
\end{response}

\section*{3. Executable manuscript commands and output}
\begin{editorcomment}
The replication script should contain all manuscript commands and print their
output; for example, \code{print(fit, threshold = 0.5)} was missing.
\end{editorcomment}
\begin{response}
We reorganized the standalone replication script to include the displayed
R input blocks, including the new comparison commands, both fit-printing calls, summaries, prediction and
reduced KIRC commands. It prints the corresponding output, regenerates the
figures and CSV tables, and records the package version and session information.
The script requires version 0.2.0 and checks the prediction table against a
separate, versioned reference. Normal mode uses the manuscript simulation and
cross-validation settings. Quick mode is a smoke test and is explicitly
separated from full numerical replication. The revised reproducibility section
and the bundle README give the commands and validation scope.
\end{response}

\section*{4. Prediction-table discrepancy and numerical correction}
\begin{editorcomment}
The binary IMR row in the expected prediction-comparison table differs from
the newly generated replication output.
\end{editorcomment}
\begin{response}
While auditing the discrepancy, we identified that the historical cross-validation
calculation used held-out responses in weighting and reused full-data selection
draws. The manuscript now evaluates prediction with
\code{cv\_imr(cv\_method = "refit")}, which fits an independent MCMC model
within each training fold and recomputes preprocessing, formula transformations
and posterior model weights from training data alone. The two post-fit modes
are retained rather than removed, and \code{"legacy"} remains the default,
because they answer a different question; the manuscript gives their
interpretation and states that they should not be labelled training-fold
refits. Hyperparameters remain fixed; data-driven tuning requires an
additional validation layer. We also corrected binary model averaging to
average model-specific probabilities, comparable-pair scoring for concordance,
and theta-pair labels when a platform occurs in four or more subgroups.

We updated the manuscript's table, interpretation and CV description. At the
manuscript settings (4000 retained draws, 1000 burn-in and ten rounds of
five-fold refit validation), the binary IMR AUC row is
\code{0.660, 0.640, 0.605, 0.605, 0.628}. All thirty entries of that table
agree with the versioned reference the replication script checks against. Two
independent full default replications produced byte-identical copies of the six
main CSV tables and eight covariate-comparison CSV tables, and identical
unrounded prediction-comparison metrics. Regression tests perturb held-out
responses and features, verify training-only transformation bases, compare
concordance against \code{survival::concordance()}, and check probability
averaging and theta labels against explicit reference calculations. The frozen
0.1.3 archive is retained for provenance, not used as the current numerical
reference.

The same audit led the previous revision to move its legacy raw-time KIRC
table into the historical record and to make no full-scale performance claim,
because that table predated the log-time preprocessing correction and used
post-fit validation. The analysis has since been rerun at the article's own
budget under corrected preprocessing, and the manuscript reports it in the
software validation section.
\end{response}

\section*{5. Software availability beyond R}
\begin{editorcomment}
The software overview focuses on R and should discuss related functionality in
other environments.
\end{editorcomment}
\begin{response}
We added a sourced cross-environment comparison table. It covers
MOFA2/\code{mofapy2} (R and Python), \code{mvlearn} (Python), similarity
network fusion (MATLAB and R), and \pkg{MultivariateStats.jl} CCA (Julia).
The discussion distinguishes latent-factor learning, correlated representations
and network fusion from IMR's availability-subgroup outcome regressions. It
states relevant input requirements and cites project documentation or the
original methods paper. This directly addresses related functionality outside
R, rather than relying on a claim that no direct IMR implementation is known.
\end{response}

\section*{6. Full-precision predictions}
\begin{editorcomment}
Predictions should retain their full values; users should choose display
precision.
\end{editorcomment}
\begin{response}
We removed internal rounding from \code{predict.imr()}. Returned probabilities
and continuous or latent-response predictions preserve numerical precision.
The manuscript example chooses \code{digits} only when printing. A regression
test verifies that returned predictions retain information beyond three decimal
places; the binary probability-averaging test separately checks numerical
correctness.
\end{response}

\section*{7. Formula/data regression interface}
\begin{editorcomment}
The package should use the usual formula/data interface, or accept a response
vector, model matrix and explicit identifiers.
\end{editorcomment}
\begin{response}
We added \code{imr.formula()}, combining a response/covariate formula and
clinical \code{data} with explicitly identified platform data. It uses
\code{model.frame()} and \code{model.matrix()} and creates the same validated
data object as the component-wise interface. Matrix-valued responses support
right-censored outcomes. Prediction retains training factor levels, contrasts
and supported transformation bases; CV learns those bases within each training
fold. Dot formulas exclude the identifier. Unsupported no-intercept and offset
formulas fail explicitly. Tests cover reordered identifiers, categorical
covariates, polynomial terms, missing variables and unseen factor levels.
The package-interface section includes the new syntax.
\end{response}

\end{document}
